\chapter{練習問題の答え}
\label{app:answers}

この付録は本文練習問題の答えを提供します。自分で解放してから確認してほしい。

\section{第1章答え}

\textbf{1-1.}ChatGPTの回答を評価してください。

ベストプラクティス：「言語モデルは次のトークン予測で学習されますが、十分に大きなモデルとさまざまなデータを使用すると、質問 - 回答パターンを暗黙的に学習します。ユーザーが質問の形で入力すると、モデルは学習データから質問の後に来た回答パターンを生成します。

\textbf{1-2.}「太陽は東から…」次の単語の確率。

人ならば、「湧き出る」を99％以上にする。理由：（1）日常的な観察、（2）言語的集団（「東から湧き出る」は固定表現）、（3）文脈の強い制約。

\textbf{1-3.}Nグラムの希少性問題。

「ホッキョクグマが踊る」次の単語をNグラムが予測するには、学習データに「ホッキョクグマが踊る」次の単語が登場する必要があります。この3グラムが学習されていない場合、確率は0になり、スムージングで「もっと」のような一般的な単語を推薦することになります。

\textbf{1-4.}トランスがRNNよりも並列処理に有利である理由。

RNNは、時点$t$の計算が時点$t-1$の隠れ状態に依存するため、順次処理必須。トランスはアテンションですべての位置を同時に処理するので、GPU並列化に最適。

\textbf{1-5.}Chinchilla 基準データスケール。

パラメータを10倍増やすと、データも約10倍増やす必要があります（20トークン/パラメータ比を維持）。つまり、10倍モデルは200倍のデータが必要。

\section{第3章 答え}

\textbf{3-1.}頻度の低いペアをマージすると？

頻繁に登場しないパターンがトークンになり、語彙が非効率的。まれな「xz」のようなペアがマージされると、実際のテキストではほとんど使用されないトークンになります。

\textbf{3-2.}Unigramペナルティ\code{-5.0}に変更。

語彙にない部分文字列のペナルティが小さくなるので、Viterbiがより積極的に分割を試みます。その結果、より小さなトークンに分解される傾向。

\textbf{3-3.}韓国語BPE、語彙100対1000。

語彙100：「こんにちは」が5つ以上のバイトトークンに分解されました。非常に非効率的。
語彙1000：「こんにちは」、「やってください」などのサブワードを学習できます。シーケンス長を大幅に短縮。

\textbf{3-4.}Viterbiの長さの上限4に縮小。

5文字以上のトークンが消えた。 「hello」のような単語が「he」、「ll」、「o」に分解され、シーケンスの長さが長くなります。

\textbf{3-5.}\code{<bos>}の必要性。

そうでなくてもモデルは動作しますが、シーケンスの開始を明示的に表示すると、モデルが「ここから新しいシーケンス」であることが認識されます。バッチ処理時に複数のシーケンスを区別するのにも便利です。

\section{第4章 答え}

\textbf{4-1.}埋め込み次元が小さすぎるか大きい場合。

小さい場合（8）：単語間の区別が難しく、意味的な類似性を学習できません。
大きな面（4096）：パラメータ爆発、メモリ不足、過適合の危険。通常256$\sim$4096を使用。

\textbf{4-2.}Sinusoidal 定数 100$\to$100.

周期が短くなり、近い位置も区別しにくくなる。位置エンコーディングの「解像度」が低くなります。

\textbf{4-3.}スケーリングを省略。

位置埋め込みがトークン埋め込みを圧倒し、モデルが位置のみに依存。意味情報の損失。

\textbf{4-4.}学習可能な位置埋め込み、max\_len 超過。

インデックス範囲外でエラーが発生しました。 RoPEのようなextrapolationが可能な方法が必要。

\textbf{4-5.}LayerNorm$\epsilon$値の変更。

$\epsilon$が大きいと正規化が弱くなり、出力分散が大きくなります。小さすぎると分母がゼロに近く、数値不安定。

\section{第5章 答え}

\textbf{5-1.}スケーリングを省略した場合のsoftmax。

$d\_k=64$のとき、内積値が$\pm 64$まで行くことができ、softmaxがワンホットに飽和。グラデーション0、学習停止。

\textbf{5-2.}因果マスク未使用。

学習時：モデルは将来のトークンを「見て」正解になるので、学習損失は低くなりますが一般化しません。
推論時：将来のトークンがないため、学習 - 推論ギャップが発生しました。パフォーマンスの急落。

\textbf{5-3.}ヘッド数 1（シングルヘッド）。

1つのアテンションパターンのみを学習。文法、意味、談話関係をしたヘッドがすべて処理しなければならないため、性能低下。

\textbf{5-4.}ヘッド数 =$d\_{model}$($d\_{head}=1$)。

各ヘッドが1次元のみ処理。アテンションの意味が消え、単純な加重になる。事実上線形変換と同じ。

\textbf{5-5.}アテンション行列メモリ、$L=100,000$。

$100000^2 \times 2$bytes (FP16) = 20 GB.シングルヘッド基準。ヘッドが32個あれば640GB。事実上不可能。

\section{第6章答え}

\textbf{6-1.} $d\_{ff} = d\_{model}$.

FFNの拡大縮小パターンが消え、非線形表現力が減少しました。より深い線形変換に似ています。

\textbf{6-2.}残差なしで12階。

backward 時グラデーションが前の階に到達できず学習不可。損失が減らない。

\textbf{6-3.}Pre-LNで\code{ln\_f}が必要です。

Pre-LNは各ブロック入力を正規化しますが、最後のブロック出力は正規化されていません。 LMヘッドに入れる前に正規化が必要です。

\textbf{6-4.}GELUの代わりにSiLU。

性能差ミミ。 SiLU（$x \cdot \sigma(x)$）もスムーズな切り替え。一部モデル（LLaMAのFFNはSwiGLUですが活性化としてSiLUを使用）が採用。

\textbf{6-5.}残差スケーリングなしで100階。

初期活性値分散が層ごとに累積されて爆発。学習の初期にNaNが発生する可能性が高い。

\section{第7章答え}

\textbf{7-1.}ウェイトタイ未使用、GPT-2小。

$V \times d = 50257 \times 768 \approx 38.6M$パラメータを追加。合計124M$\to$163M。

\textbf{7-2.} Warmup 0.

学習初期グラデーションが不安定で発散の可能性。特にPost-LNでは深刻。 Pre-LNはそれほど敏感ではありませんが、まだお勧めします。

\textbf{7-3.} Label smoothing 0.3.

強いスムージングで、モデルは正解に対する確信が少ない。一般化は向上しますが、学習損失が高く維持。通常0.1が適正。

\textbf{7-4.}グラデーションクリッピングなし。

グラデーション爆発時の重みが大きく変化し、モデルが崩壊。特に大きな学習率で危険。

\textbf{7-5.}PPL 1.0 可能ですか？

理論的に不可能。自然言語は本質的に不確実性があり、次のトークンが100\%予測不可。学習データに完全に覚えている場合にのみ、特定の文でPPL$\approx$1可能ですが、一般的ではありません。

\section{第8章 答え}

\textbf{8-1.}Greedy 繰り返しループの原因。

分布が「store」の後に「to buy」を高く、「to buy」の後に「some store」を高く予測するループが形成されると、greedyはそこから抜け出すことができません。確率の高い経路だけに従うからだ。

\textbf{8-2.} Temperature $\to$ 0.

$T=0$なら softmax が argmax になり greedy と同じ。$T=0.01$の場合、ほぼgreedy。分布が鋭すぎて多様性 0.

\textbf{8-3.}Top-K後Top-P。

Top-Kで候補をK個に減らした後、その中から累積確率Pまで選択。より強力なフィルタリング。品質は向上しますが、多様性の減少。

\textbf{8-4.}チャットボットに適したサンプリング。

Top-P (P=0.9、T=0.7$\sim$0.9)。多様性と一貫性のバランス。 Greedyは単調で、純粋なサンプリングは一貫性不足。

\textbf{8-5.}Speculative decodingが速い理由。

小さなモデルがK個のトークンをすばやく生成すると、大きなモデルがそのK個を一度に検証（並列処理）。当たるとK個を1ステップに採用するので、大型モデルのステップ数が1/Kに減少。
